\documentclass[12pt]{report}
\usepackage[T1]{fontenc}
\usepackage[brazilian]{babel}
\usepackage{csquotes}
\usepackage[
  perfil=academico,
  tipo-trabalho=dissertacao,
  impressao=frente-e-verso,
  citacoes=autor-data
]{abntex3}

\addbibresource{referencias-exemplo.bib}

\ABNTEXsetup{
  titulo={Perfil acadêmico no abnTeX3 Community},
  subtitulo={exemplo integrado do Marco 7},
  autoria={Maria da Silva},
  instituicao={Universidade Exemplo},
  natureza={Dissertação},
  objetivo={apresentada para obtenção do título de Mestre},
  area={Documentação},
  orientacao={Prof. Dr. João Souza},
  local={Recife},
  data={2026}
}

\begin{document}

\ABNTEXacademicoexterna

\ABNTEXacademicopretextual
\ABNTEXfolhaderosto

% Elemento opcional: substitua o conteúdo pelos dados preparados pelo
% profissional responsável ou remova o comando inteiro.
\ABNTEXfichacatalografica{%
  Dados internacionais de catalogação na publicação.
}

\ABNTEXfolhadeaprovacao
  {29 de julho de 2026}
  {%
    \ABNTEXmembrobanca
      {Prof. Dr. João Souza}{Orientador}{Universidade Exemplo}
    \ABNTEXmembrobanca
      {Profa. Dra. Ana Santos}{Examinadora}{Universidade Exemplo}
  }

\ABNTEXresumo
  {Este exemplo reúne os módulos do pacote em um fluxo completo de trabalho
   acadêmico, da capa às referências.}
  {normalização, trabalho acadêmico, LaTeX}

\ABNTEXresumoemlinguaestrangeira
  {This example combines the package modules into a complete academic work
   flow, from the cover to the references.}
  {standardization, academic work, LaTeX}

\ABNTEXsumario

\ABNTEXacademicotextual

\chapter{Introdução}

O perfil preserva a divisão de responsabilidades: apresentação documental
fica no abnTeX3 Community, enquanto citações e referências são processadas
pelo \texttt{biblatex-abnt} e pelo \texttt{biber}
\parencite{associacao2025}.

\chapter{Desenvolvimento}

Os quatro tipos de trabalho compartilham o mesmo fluxo. A chave de tipo
identifica a configuração, e os metadados institucionais permanecem
explícitos.

\ABNTEXacademicopostextual
\ABNTEXbibliografia

\end{document}
